\section{Introduction}
\label{sec:intro}

\begin{figure}[t]
    \centering
    \includegraphics[width=\columnwidth]{figs/head.pdf}
    \caption{a) Conventional CL learns distinct task-specific heads, while CL with vision-language models can predict both learned tasks and out-of-distribution tasks. b) Accuracy (\%) changes during CL of four datasets on 11 datasets. Our method is superior to others in preventing the forgetting of both zero-shot transfer ability and new knowledge.}
    \label{fig:head}
\end{figure}

% Background: Continual learning
Most deep learning models can access all the data during training~\cite{krizhevsky2017imagenet,he2016deep,dosovitskiy2020image}.
% are trained with all the data available
% \kai{say these methods can access all the data during the whole training}
If we want to expand a model's knowledge, such as learning a newly found animal species~\cite{Parkhi2012CatsAD}, we can re-train the model by adding new classes to the training dataset. However, re-training a large model is costly.
In contrast, continual learning (CL)~\cite{li2017learning,rebuffi2017icarl} incrementally learns task one after another.
% acquires incremental knowledge of various tasks over time. 
% \kai{say the training data of CL is task by task in an order or other}
It can reduce this cost by only learning the new data, thereby presenting itself as an efficient alternative to conventional learning methods~\cite{qin2023infobatch,liu2023dream,zhou2023dataset}.
% \kai{this sentence should say more clearly}
Nonetheless, a model tends to forget previous information catastrophically when learning new tasks~\cite{rebuffi2017icarl,li2017learning,wu2019large}. 
% \kai{maybe here is access to earlier data?}
The ``catastrophic forgetting'' phenomenon is a great challenge for CL.
% \kai{it is better to show some results,such as ref to fig 1?}

% Vision-Language Model intro; CL benefit
Recently, vision-language models have shown powerful zero-shot transfer ability~\cite{radford2021learning, jia2021scaling, Li2022BLIPBL}. They can give zero-shot predictions without any training examples of a task.
% \kai{we can say directly the powerful zero-shot ability}
However, the performance on some tasks is poor due to insufficient relevant image-text pairs in the pre-training datasets. For example, it is difficult for CLIP~\cite{radford2021learning} to distinguish among digital numbers, with an accuracy on MNIST~\cite{deng2012mnist} below 60\%
% between an Airbus and a Boeing airplane~\cite{maji13fine-grained}, and the accuracy on MNIST~\cite{deng2012mnist} is below 60\%, 
much lower than a naively trained CNN~\cite{lecun1998gradient}. If we want to widen the knowledge in the vision-language model by re-training, the computational cost is too large (e.g., CLIP is pretrained on 400 million image-text pairs). Fine-tuning downstream tasks achieves high performance, but one model for a task takes much memory, and the model is not reusable.
% Besides, it takes minutes to fine-tune. The latency is unacceptable to fine-tune when needed.
% the learned model cannot generalize outside the fine-tuned dataset. \kai{here, it needs to consider to compare the cost between faster fine-tune and directly train on some datasets. Nobody would like to rerun CLIP, but some guys may concern whether we can faster fine to achieve sota results based on clip}
% Linear-probing~\cite{radford2021learning} and prompt learning~\cite{zhou2022learning,zhou2022conditional} keep the backbone parameters unchanged. However, the task id must be known in advance to select the corresponding classification heads and prompts. 
Prompt learning~\cite{zhou2022learning,zhou2022conditional} keeps the backbone parameters unchanged. However, it is only effective with limited training data due to a limited prompt length~\cite{zhou2022learning,jin2021good}.
% \kai{provide briefly analysis why they fail to large-scale dataset}
In contrast, continual learning makes learning new knowledge a lifelong process for the vision-language model.
% \kai{we can not claim this, as cl is just a way to integrate novel knowledge into vl model, change the current claim to avoid confusion.}
The continually learned model can handle any image-text input and can be further used for downstream tasks~\cite{ding2022don,thengane2022clip}. 
% \kai{any references to support this sentence?}

% previous methods fail
We find that existing CL methods hardly prevent the forgetting phenomenon for zero-shot transfer ability in continual learning of a pre-trained vision-language model. As shown in \cref{fig:head} (a), the CL with a pre-trained vision-language model differs from the traditional one. Besides forgetting previously learned task knowledge, the CLIP-based CL suffers from forgetting pre-training knowledge, namely a degradation of zero-shot transfer ability. For the replay-based CL methods~\cite{rebuffi2017icarl,shin2017continual,lopez2017gradient,isele2018selective,lavda2018continual,prabhu2020gdumb}, the dataset during pre-training may be private and inaccessible during fine-tuning. For distillation-based CL methods~\cite{li2017learning,dhar2019learning,douillard2020podnet,ding2022don}, they do not lay enough emphasis on the pre-trained model. On the one hand, a large model state change hinders tasks thereafter from using high-quality feature representations. On the other hand, it significantly degrades zero-shot performance on unseen datasets.

% our novelty: our method
Our method ZSCL protects the Zero-Shot transfer ability during Continual Learning. We view the knowledge stored in the pre-training model from two perspectives: a well-learned feature space and a good value in the parameter space. In feature space, we re-design previous distillation loss~\cite{hinton2015distilling, rebuffi2017icarl} with different loss styles, teacher models, and data sources. We find the original CLIP model, as opposed to the newly acquired model, is the best option for the teacher model. Instead of using data collected from previous tasks~\cite{rebuffi2017icarl} or current task~\cite{hinton2015distilling}, we find a reference dataset with diverse semantics (e.g., images sampled from ImageNet) is a good option for distillation loss. The reference images need not be labeled or matched with the text. Preserving the relative similarity between reference images and texts makes the feature space deviate little from the original. In the parameter space, WiSE-FT~\cite{wortsman2022robust} proposes interpolating the initial and fine-tuned model for better performance. Inspired by this, we ensemble the weights throughout continuous training to prevent a significant shift from the initial CLIP, which can be seen as interpolating models of different zero-shot transfer and downstream task performance tradeoffs. The weight ensemble method is more stable and not sensitive to hyper-parameters.

% experiment performance
To better evaluate our method, we propose a new benchmark Multi-domain Task Incremental Learning (MTIL). Previous CL tasks are crafted by separating classes in one dataset~\cite{douillard2022dytox,yan2021dynamically,zhu2021prototype}, where the images and classes are within a single domain. In contrast, MTIL consists of data from different sources requiring different expert knowledge. It comprises 11 tasks ranging from animal species to aircraft series recognition. As displayed in \cref{fig:head} (b), when sequentially training CLIP on 11 datasets, the drop in the performance of task $i$ after training task $i$ is the traditional forgetting phenomenon. The degradation in the accuracy compared to the original zero-shot one before training task $i$ represents the forgetting in zero-shot transfer ability. Our method better protects the zero-shot transfer ability and preserves the learned knowledge. We outperform previous methods in both conventional class-incremental learning and MTIL settings.
In \cref{fig:head} (b), 

% our contribution
To summarize, our contributions are as follows:
\begin{itemize}
    \item We investigate continual training with the vision-language model and demonstrate the importance of preserving zero-shot transfer ability. A more challenging benchmark MTIL is proposed to evaluate CL methods where the tasks come from distinct domains.
    \item We propose a novel method ZSCL to mitigate the catastrophic forgetting problem in continual learning of the vision-language model by distillation in the feature space and weight ensemble in the parameter space. 
    \item The proposed ZSCL outperforms all state-of-the-art methods across multiple benchmark datasets. On 10 steps CL of CIFAR100 and TinyImageNet, our method outperforms the best of previous ones by 7.7\% and 6.0\% for the Last accuracy. On MTCL, ZSCL outperforms others by 10.9\% on Transfer and 9.7\% on Avg. scores.
\end{itemize}